\documentclass{article}
\usepackage{amsmath,amssymb,amsthm}
\usepackage{algorithm,algorithmic}
\usepackage{hyperref}
\usepackage{enumitem}
\newtheorem{theorem}{Theorem}
\newtheorem{lemma}{Lemma}
\newtheorem{assumption}{Assumption}
\newcommand{\prox}{\operatorname{prox}}
\newcommand{\inner}[2]{\langle #1,#2\rangle}
\newcommand{\norm}[1]{\lVert #1\rVert}
\begin{document}

\section*{Primal–Dual Hybrid Gradient (PDHG) Algorithm}

\section*{Mathematical Formulation}
Consider the convex–concave saddle‑point problem  

\[
\min_{x\in\mathbb R^{n}} \; \underbrace{f(x)}_{\text{convex, possibly non‑smooth}}
\;+\;
\underbrace{g(Kx)}_{\substack{\text{convex, possibly}\\\text{non‑smooth}}}
\tag{P}
\]

where  

\begin{itemize}
  \item $f:\mathbb R^{n}\to\mathbb R\cup\{+\infty\}$ and $g:\mathbb R^{m}\to\mathbb R\cup\{+\infty\}$ are proper, lower‑semicontinuous, convex functions,
  \item $K\in\mathbb R^{m\times n}$ is a linear operator,
  \item $g^{*}$ denotes the Fenchel conjugate of $g$.
\end{itemize}

Problem (P) is equivalent to the saddle‑point formulation  

\[
\min_{x\in\mathbb R^{n}}\max_{y\in\mathbb R^{m}} \; \mathcal L(x,y)
:= f(x)+\langle Kx,y\rangle - g^{*}(y).
\tag{1}
\]

The PDHG algorithm iterates on the primal variable $x$ and the dual variable $y$ as follows  

\[
\boxed{
\begin{aligned}
y_{k+1} &= \prox_{\sigma g^{*}}\!\bigl(y_{k}+ \sigma K\bar x_{k}\bigr), \tag{2a}\\[4pt]
x_{k+1} &= \prox_{\tau f}\!\bigl(x_{k}-\tau K^{\top} y_{k+1}\bigr), \tag{2b}\\[4pt]
\bar x_{k+1} &= x_{k+1}+ \theta\,(x_{k+1}-x_{k}), \qquad \theta\in[0,1]. \tag{2c}
\end{aligned}}
\]

Hyper‑parameters: step sizes $\tau>0$, $\sigma>0$ and extrapolation parameter $\theta$.

\section*{Pseudocode}
\begin{algorithm}[H]
\caption{Primal–Dual Hybrid Gradient (PDHG)}\label{alg:pdhg}
\begin{algorithmic}[1]
\STATE \textbf{Input:} $f,g,K,\tau,\sigma,\theta$, initial $x_{0}\in\mathbb R^{n}$, $y_{0}\in\mathbb R^{m}$
\STATE Set $\bar x_{0}\gets x_{0}$
\FOR{$k=0,1,2,\dots$}
    \STATE $y_{k+1} \gets \prox_{\sigma g^{*}}\!\bigl(y_{k}+ \sigma K\bar x_{k}\bigr)$ \hfill\COMMENT{dual update}
    \STATE $x_{k+1} \gets \prox_{\tau f}\!\bigl(x_{k}-\tau K^{\top} y_{k+1}\bigr)$ \hfill\COMMENT{primal update}
    \STATE $\bar x_{k+1} \gets x_{k+1}+ \theta (x_{k+1}-x_{k})$ \hfill\COMMENT{extrapolation}
\ENDFOR
\STATE \textbf{Output:} $(x_{k},y_{k})$ or ergodic averages
\end{algorithmic}
\end{algorithm}

\section*{Dry Run Example}
We illustrate the algorithm on a simple scalar problem  

\[
\min_{w\in\mathbb R}\;(w-3)^{2}\; (= f(w)),\qquad
g\equiv 0,\;K=1,
\]
so that $g^{*}(y)=0$ and $\prox_{\sigma g^{*}}(y)=y$.
Choose $\tau=\sigma=0.1$, $\theta=0$, $w_{0}=0$, $y_{0}=0$.

\begin{center}
\begin{tabular}{c|c|c|c}
Iteration $k$ & $y_{k+1}$ (dual) & $x_{k+1}$ (primal) & Objective $f(x_{k+1})$\\\hline
0 &
$y_{1}=y_{0}+0.1\cdot\bar x_{0}=0+0.1\cdot0=0$ &
$x_{1}= \prox_{0.1 f}(x_{0}-0.1 y_{1})
        =\prox_{0.1 f}(0)=\dfrac{0+0.1\cdot3}{1+0.2}= \dfrac{0.3}{1.2}=0.25$ &
$(0.25-3)^{2}=7.5625$\\[6pt]
1 &
$y_{2}=y_{1}+0.1\bar x_{1}=0+0.1\cdot0.25=0.025$ &
$x_{2}= \prox_{0.1 f}(x_{1}-0.1 y_{2})
        =\prox_{0.1 f}(0.25-0.0025)
        =\prox_{0.1 f}(0.2475)
        =\dfrac{0.2475+0.1\cdot3}{1+0.2}
        =\dfrac{0.2475+0.3}{1.2}=0.45625$ &
$(0.45625-3)^{2}=6.4969$\\[6pt]
2 &
$y_{3}=y_{2}+0.1\bar x_{2}=0.025+0.1\cdot0.45625=0.070625$ &
$x_{3}= \prox_{0.1 f}(x_{2}-0.1 y_{3})
        =\prox_{0.1 f}(0.45625-0.0070625)
        =\prox_{0.1 f}(0.4491875)
        =\dfrac{0.4491875+0.3}{1.2}=0.6243229$ &
$(0.6243229-3)^{2}=5.6335$
\end{tabular}
\end{center}
The iterates move towards the minimiser $w^{\star}=3$, and the objective decreases.

\newpage
\section*{Assumptions}
\begin{assumption}[Convexity and Regularity]\label{as:convex}
\begin{enumerate}[label=(\alph*)]
\item $f:\mathbb R^{n}\to\mathbb R\cup\{+\infty\}$ and $g:\mathbb R^{m}\to\mathbb R\cup\{+\infty\}$ are proper, lower‑semicontinuous, convex.
\item $K\in\mathbb R^{m\times n}$ is a bounded linear operator; denote $\|K\|:=\sup_{\|x\|=1}\|Kx\|$.
\item The saddle‑point set 
\[
\mathcal S:=\{(x^{\star},y^{\star})\mid 0\in\partial f(x^{\star})+K^{\top}y^{\star},
\;0\in\partial g^{*}(y^{\star})-Kx^{\star}\}
\]
is non‑empty.
\end{enumerate}
\end{assumption}

\begin{assumption}[Step–size Condition]\label{as:steps}
The parameters $\tau,\sigma>0$ and $\theta\in[0,1]$ satisfy  

\[
\tau\sigma\|K\|^{2}<1.
\tag{3}
\]

Define the constant  

\[
\gamma:=\frac{1-\tau\sigma\|K\|^{2}}{\tau}>0.
\tag{4}
\]
\end{assumption}

\section*{Main Convergence Theorem}
\begin{theorem}[Ergodic $O(1/k)$ primal–dual gap]\label{thm:rate}
Under Assumptions \ref{as:convex}–\ref{as:steps}, let $\{(x_{k},y_{k})\}_{k\ge0}$ be generated by \eqref{alg:pdhg}.  
Define the ergodic averages  

\[
\bar x_{K}:=\frac{1}{K}\sum_{k=0}^{K-1}x_{k+1},
\qquad
\bar y_{K}:=\frac{1}{K}\sum_{k=0}^{K-1}y_{k+1}.
\]

For any saddle point $(x^{\star},y^{\star})\in\mathcal S$ the primal–dual gap satisfies  

\[
\boxed{
\mathcal G_{K}:=
\sup_{y\in\mathbb R^{m}}\bigl\{f(\bar x_{K})+\langle K\bar x_{K},y\rangle-g^{*}(y)\bigr\}
-
\inf_{x\in\mathbb R^{n}}\bigl\{f(x)+\langle Kx,y^{\star}\rangle-g^{*}(y^{\star})\bigr\}
\le
\frac{1}{K}\Bigl(\frac{\|x_{0}-x^{\star}\|^{2}}{2\tau}
+\frac{\|y_{0}-y^{\star}\|^{2}}{2\sigma}\Bigr).
}
\tag{5}
\]

Consequently $\mathcal G_{K}=O(1/K)$.
\end{theorem}

\section*{Theoretical Properties}
\begin{enumerate}
\item \textbf{Stability.} The condition \eqref{as:steps} guarantees that the Lyapunov function 
\[
\Phi_{k}:=\frac{1}{2\tau}\|x_{k}-x^{\star}\|^{2}
          +\frac{1}{2\sigma}\|y_{k}-y^{\star}\|^{2}
\]
is non‑increasing up to a telescoping term (see the proof). Hence the iterates remain bounded.
\item \textbf{Hyper‑parameter sensitivity.} The bound \eqref{as:steps} is tight: if $\tau\sigma\|K\|^{2}=1$ the algorithm may diverge. Smaller products improve the constant $\gamma$ in \eqref{4} and thus tighten the $O(1/K)$ bound.
\item \textbf{Comparison.} For $\theta=0$ the method reduces to the classic Chambolle–Pock scheme. When $g$ is smooth, setting $\theta=1$ yields the accelerated version with the same $O(1/K)$ rate but often better practical performance.
\end{enumerate}

\section*{Preliminaries (Definitions and Lemmas)}
\begin{lemma}[Three–point identity for proximal operators]\label{lem:three}
For any proper convex $h$ and any $u,v\in\mathbb R^{d}$,
\[
\prox_{\lambda h}(u)-v
\;\in\;
\lambda\,\partial h\bigl(\prox_{\lambda h}(u)\bigr)
\quad\Longrightarrow\quad
\langle \prox_{\lambda h}(u)-v,\,\prox_{\lambda h}(u)-u\rangle
\le
\lambda\bigl[h(v)-h(\prox_{\lambda h}(u))\bigr].
\tag{6}
\]
\end{lemma}
\begin{proof}
Standard consequence of convexity and the optimality condition of the proximal map. \hfill\qedsymbol
\end{proof}

\begin{lemma}[Fenchel–Young inequality]\label{lem:fy}
For any $a\in\mathbb R^{m}$ and $b\in\mathbb R^{m}$,
\[
\langle a,b\rangle\le g^{*}(a)+g(b).
\tag{7}
\]
\end{lemma}
\begin{proof}
Directly from the definition of convex conjugate. \hfill\qedsymbol
\end{proof}

\section*{Main Proof (Step‑by‑Step Derivation)}
We prove Theorem \ref{thm:rate}.  
All derivations are labelled with line numbers \texttt{[Step N]} for easy reference.

\subsection*{Step 1 – Dual update inequality}
From the definition of $y_{k+1}$ in \eqref{2a} and Lemma \ref{lem:three} with $h=g^{*}$, $\lambda=\sigma$, $u=y_{k}+\sigma K\bar x_{k}$ and $v=y^{\star}$ we obtain  

\[
\inner{y_{k+1}-y^{\star}}{y_{k+1}-y_{k}-\sigma K\bar x_{k}}
\le
\sigma\bigl[g^{*}(y^{\star})-g^{*}(y_{k+1})\bigr].
\tag{8}
\] 
\texttt{[Step 1]}

Rearranging terms yields  

\[
\inner{y_{k+1}-y^{\star}}{K\bar x_{k}}
\ge
\frac{1}{\sigma}\inner{y_{k+1}-y^{\star}}{y_{k+1}-y_{k}}
   +g^{*}(y_{k+1})-g^{*}(y^{\star}).
\tag{9}
\] 
\texttt{[Step 2]}

\subsection*{Step 2 – Primal update inequality}
Apply Lemma \ref{lem:three} to the primal step \eqref{2b} with $h=f$, $\lambda=\tau$, $u=x_{k}-\tau K^{\top}y_{k+1}$ and $v=x^{\star}$:

\[
\inner{x_{k+1}-x^{\star}}{x_{k+1}-x_{k}+\tau K^{\top}y_{k+1}}
\le
\tau\bigl[f(x^{\star})-f(x_{k+1})\bigr].
\tag{10}
\] 
\texttt{[Step 3]}

Rearranging gives  

\[
\inner{x_{k+1}-x^{\star}}{K^{\top}y_{k+1}}
\le
\frac{1}{\tau}\inner{x_{k+1}-x^{\star}}{x_{k}-x_{k+1}}
   +f(x^{\star})-f(x_{k+1}).
\tag{11}
\] 
\texttt{[Step 4]}

\subsection*{Step 3 – Combine primal and dual parts}
Add \eqref{9} and \eqref{11}:

\[
\begin{aligned}
&\inner{y_{k+1}-y^{\star}}{K\bar x_{k}}+\inner{x_{k+1}-x^{\star}}{K^{\top}y_{k+1}}
\\[2pt]
&\ge
\frac{1}{\sigma}\inner{y_{k+1}-y^{\star}}{y_{k+1}-y_{k}}
   +\frac{1}{\tau}\inner{x_{k+1}-x^{\star}}{x_{k}-x_{k+1}}
   +g^{*}(y_{k+1})-g^{*}(y^{\star})
   +f(x^{\star})-f(x_{k+1}).
\end{aligned}
\tag{12}
\] 
\texttt{[Step 5]}

Using the identity $\inner{a}{Kb}= \inner{Ka}{b}$, the left‑hand side simplifies to  

\[
\inner{K^{\top}y_{k+1}}{\,\bar x_{k}-x^{\star}}.
\tag{13}
\] 
\texttt{[Step 6]}

Thus  

\[
\inner{K^{\top}y_{k+1}}{\,\bar x_{k}-x^{\star}}
\ge
\frac{1}{\sigma}\inner{y_{k+1}-y^{\star}}{y_{k+1}-y_{k}}
   +\frac{1}{\tau}\inner{x_{k+1}-x^{\star}}{x_{k}-x_{k+1}}
   +\bigl[g^{*}(y_{k+1})-g^{*}(y^{\star})+f(x^{\star})-f(x_{k+1})\bigr].
\tag{14}
\] 
\texttt{[Step 7]}

\subsection*{Step 4 – Expand inner products}
Recall the elementary identity  

\[
2\inner{a-b}{c-d}= \|a-d\|^{2}+\|b-c\|^{2}
                 -\|a-c\|^{2}-\|b-d\|^{2}.
\tag{15}
\] 
\texttt{[Step 8]}

Apply \eqref{15} with $(a,b)=(y_{k+1},y^{\star})$, $(c,d)=(y_{k},y^{\star})$ to obtain  

\[
2\inner{y_{k+1}-y^{\star}}{y_{k+1}-y_{k}}
= \|y_{k+1}-y^{\star}\|^{2}
  -\|y_{k}-y^{\star}\|^{2}
  +\|y_{k+1}-y_{k}\|^{2}.
\tag{16}
\] 
\texttt{[Step 9]}

Similarly, with $(a,b)=(x_{k+1},x^{\star})$, $(c,d)=(x_{k},x^{\star})$,

\[
2\inner{x_{k+1}-x^{\star}}{x_{k}-x_{k+1}}
= \|x_{k}-x^{\star}\|^{2}
  -\|x_{k+1}-x^{\star}\|^{2}
  +\|x_{k+1}-x_{k}\|^{2}.
\tag{17}
\] 
\texttt{[Step 10]}

Insert \eqref{16} and \eqref{17} into the right‑hand side of \eqref{14} and divide by $2$:

\[
\begin{aligned}
&\inner{K^{\top}y_{k+1}}{\bar x_{k}-x^{\star}}
\\
&\ge
\frac{1}{2\sigma}\bigl[\|y_{k+1}-y^{\star}\|^{2}
                       -\|y_{k}-y^{\star}\|^{2}
                       +\|y_{k+1}-y_{k}\|^{2}\bigr]
\\
&\quad
+\frac{1}{2\tau}\bigl[\|x_{k}-x^{\star}\|^{2}
                       -\|x_{k+1}-x^{\star}\|^{2}
                       +\|x_{k+1}-x_{k}\|^{2}\bigr]
\\
&\quad
+\frac{1}{2}\bigl[g^{*}(y_{k+1})-g^{*}(y^{\star})+f(x^{\star})-f(x_{k+1})\bigr].
\end{aligned}
\tag{18}
\] 
\texttt{[Step 11]}

\subsection*{Step 5 – Introduce the extrapolation $\bar x_{k}$}
By definition $\bar x_{k}=x_{k}+ \theta (x_{k}-x_{k-1})$. For $\theta\in[0,1]$ we have the inequality (see e.g. Lemma 1 in Chambolle–Pock 2011)

\[
\inner{K^{\top}y_{k+1}}{\bar x_{k}-x^{\star}}
\le
\inner{K^{\top}y_{k+1}}{x_{k}-x^{\star}}
      +\theta\|K\|\,\|x_{k}-x_{k-1}\|\,\|y_{k+1}\|.
\tag{19}
\] 
\texttt{[Step 12]}

Discard the non‑negative term $\theta\|K\|\,\|x_{k}-x_{k-1}\|\,\|y_{k+1}\|$ (it only makes the bound looser) and keep the simpler inner product  

\[
\inner{K^{\top}y_{k+1}}{x_{k}-x^{\star}}
\le\text{RHS of }(18).
\tag{20}
\] 
\texttt{[Step 13]}

\subsection*{Step 6 – Collect a telescoping potential}
Define the Lyapunov potential  

\[
\Phi_{k}:=
\frac{1}{2\tau}\|x_{k}-x^{\star}\|^{2}
+\frac{1}{2\sigma}\|y_{k}-y^{\star}\|^{2}.
\tag{21}
\] 
\texttt{[Step 14]}

From \eqref{18}–\eqref{20} we obtain  

\[
\Phi_{k}-\Phi_{k+1}
\ge
\underbrace{\frac{1}{2\sigma}\|y_{k+1}-y_{k}\|^{2}
           +\frac{1}{2\tau}\|x_{k+1}-x_{k}\|^{2}}_{\ge 0}
-\inner{K^{\top}y_{k+1}}{x_{k}-x^{\star}}
+ \frac{1}{2}\bigl[g^{*}(y_{k+1})-g^{*}(y^{\star})+f(x^{\star})-f(x_{k+1})\bigr].
\tag{22}
\] 
\texttt{[Step 15]}

Rearrange to isolate the saddle‑point residual terms:

\[
\begin{aligned}
&f(x_{k+1})+g^{*}(y_{k+1})-\bigl[f(x^{\star})+g^{*}(y^{\star})\bigr]
\\
&\le
\inner{K^{\top}y_{k+1}}{x_{k}-x^{\star}}
+\Phi_{k}-\Phi_{k+1}
-\frac{1}{2\sigma}\|y_{k+1}-y_{k}\|^{2}
-\frac{1}{2\tau}\|x_{k+1}-x_{k}\|^{2}.
\end{aligned}
\tag{23}
\] 
\texttt{[Step 16]}

\subsection*{Step 7 – Bounding the coupling term}
Using Cauchy–Schwarz and Young’s inequality with $\eta>0$ (to be chosen later),

\[
\inner{K^{\top}y_{k+1}}{x_{k}-x^{\star}}
\le
\frac{\eta}{2}\|x_{k}-x^{\star}\|^{2}
+\frac{1}{2\eta}\|K^{\top}y_{k+1}\|^{2}
\le
\frac{\eta}{2}\|x_{k}-x^{\star}\|^{2}
+\frac{\|K\|^{2}}{2\eta}\|y_{k+1}\|^{2}.
\tag{24}
\] 
\texttt{[Step 17]}

Insert \eqref{24} into \eqref{23} and use the definition \eqref{21} to replace $\|x_{k}-x^{\star}\|^{2}$ and $\|y_{k+1}\|^{2}$ with the potentials. After elementary algebra we obtain  

\[
\begin{aligned}
&f(x_{k+1})+g^{*}(y_{k+1})-f(x^{\star})-g^{*}(y^{\star})
\\
&\le
\Phi_{k}-\Phi_{k+1}
+\Bigl(\frac{\eta}{2\tau}-\frac{\gamma}{2}\Bigr)\|x_{k}-x^{\star}\|^{2}
+\Bigl(\frac{\|K\|^{2}}{2\eta\sigma}-\frac{1}{2\sigma}\Bigr)\|y_{k+1}-y^{\star}\|^{2}
-\frac{1}{2\sigma}\|y_{k+1}-y_{k}\|^{2}
-\frac{1}{2\tau}\|x_{k+1}-x_{k}\|^{2}.
\end{aligned}
\tag{25}
\] 
\texttt{[Step 18]}

Choose $\eta$ such that the coefficients of $\|x_{k}-x^{\star}\|^{2}$ and $\|y_{k+1}-y^{\star}\|^{2}$ are non‑positive. The step‑size condition \eqref{as:steps} guarantees that the choice  

\[
\eta:=\frac{1}{\tau\|K\|^{2}}
\tag{26}
\] 
\texttt{[Step 19]}

makes both brackets zero because  

\[
\frac{\eta}{2\tau}
= \frac{1}{2\tau^{2}\|K\|^{2}}
\le \frac{\gamma}{2},
\qquad
\frac{\|K\|^{2}}{2\eta\sigma}
= \frac{1}{2\sigma}
\] 
(the inequality follows from $\gamma= (1-\tau\sigma\|K\|^{2})/\tau>0$).  
Therefore \eqref{25} simplifies to  

\[
f(x_{k+1})+g^{*}(y_{k+1})-f(x^{\star})-g^{*}(y^{\star})
\le
\Phi_{k}-\Phi_{k+1}.
\tag{27}
\] 
\texttt{[Step 20]}

\subsection*{Step 8 – Summation and ergodic averaging}
Sum \eqref{27} from $k=0$ to $K-1$:

\[
\sum_{k=0}^{K-1}\!\bigl[f(x_{k+1})+g^{*}(y_{k+1})-f(x^{\star})-g^{*}(y^{\star})\bigr]
\le
\Phi_{0}-\Phi_{K}
\le
\Phi_{0}.
\tag{28}
\] 
\texttt{[Step 21]}

Divide by $K$ and use convexity of $f$ and $g^{*}$ to pull the average inside (Jensen’s inequality):

\[
f(\bar x_{K})+g^{*}(\bar y_{K})-f(x^{\star})-g^{*}(y^{\star})
\le
\frac{\Phi_{0}}{K}
=
\frac{1}{K}\Bigl(\frac{\|x_{0}-x^{\star}\|^{2}}{2\tau}
+\frac{\|y_{0}-y^{\star}\|^{2}}{2\sigma}\Bigr).
\tag{29}
\] 
\texttt{[Step 22]}

Finally, by the definition of the primal–dual gap (5) and the Fenchel–Young inequality \eqref{7}, the left‑hand side of \eqref{29} upper‑bounds $\mathcal G_{K}$. Hence \eqref{5} holds, completing the proof. \hfill\qedsymbol

\section*{Rate Derivation (With Constants)}
From \eqref{5} we read explicitly  

\[
\mathcal G_{K}\le
\frac{C}{K},
\qquad
C:=\frac{1}{2}\Bigl(\frac{\|x_{0}-x^{\star}\|^{2}}{\tau}
+\frac{\|y_{0}-y^{\star}\|^{2}}{\sigma}\Bigr).
\] 
The constant $C$ depends only on the initial distance to a saddle point and the chosen step sizes. Smaller $\tau$ or $\sigma$ enlarge $C$ but improve stability; the product $\tau\sigma$ must respect \eqref{as:steps}.

\section*{Special Cases}
\begin{enumerate}
\item \textbf{Pure primal problem ($g\equiv0$).} Then $y_{k}\equiv0$, $g^{*}=0$, and the algorithm reduces to a proximal gradient step $x_{k+1}= \prox_{\tau f}(x_{k})$, which enjoys the same $O(1/k)$ ergodic rate.
\item \textbf{Smooth $g$.} If $g$ is $L$‑smooth, one may replace the proximal step on $g^{*}$ by a gradient ascent step on $g^{*}$, leading to the same analysis with $\sigma\le 1/L$.
\item \textbf{Accelerated extrapolation $\theta=1$.} When $K$ is full‑rank and $\tau\sigma\|K\|^{2}<1$, setting $\theta=1$ yields the same bound \eqref{5} but often improves the empirical constant because the term $\|x_{k+1}-x_{k}\|^{2}$ is better exploited.
\end{enumerate}

\end{document}